\chapter{Output, Diagnostics, and Results}
\label{chap:output-results}

FEMaster writes analysis results in its own result formats and provides several diagnostic commands for inspecting matrices, constraints, and the assembled model. These diagnostics do not change the physical problem; they are intended for verification, comparison, and troubleshooting.

The native \file{.res} format preserves FEMaster Field domains and is the most complete textual result source. FEMR provides a seekable, compressed binary representation for efficient programmatic access. FRD output provides a CalculiX/CGX-compatible visualization format for supported nodal quantities.

Abaqus output-request cards such as \cmd{NODE FILE} are accepted where documented so that otherwise supported Abaqus decks do not fail solely because they contain common output requests. These cards currently do not control FEMaster's result writers.

\keywordpage{OVERVIEW}

\cmd{OVERVIEW} prints a diagnostic summary of the current model without changing the analysis.

The overview is intended to answer basic model-verification questions before numerical results are interpreted: which Parts and Instances exist, how many nodes and elements are present, which named sets and surfaces are available, which materials and Sections have been defined, and which loads, supports, and constraints are present.

In models with a Part/Instance hierarchy or many named regions, the overview exposes missing Instances, empty sets, unexpected entity counts, and omitted property assignments independently of the linear solver.

A structurally plausible overview establishes only the reported model organization. It does not establish the physical consistency of material units, element orientations, shell normals, surface sides, load directions, or boundary conditions.

\exampleheading
\begin{dialectexamples}
\begin{femasterdeck}
*OVERVIEW
\end{femasterdeck}
\begin{noabaqus}
No Abaqus input keyword maps directly to FEMaster's native \cmd{OVERVIEW} diagnostic.
\end{noabaqus}
\end{dialectexamples}

\keywordpage{REQUESTSTIFFNESS}

\cmd{REQUESTSTIFFNESS} requests stiffness-matrix output from supported native analyses. It is available for linear static, nonlinear static, and linear buckling loadcases. \key{FILE} optionally specifies the output filename; otherwise FEMaster generates a deterministic default name.

For linear static analysis, the exported matrix exposes symmetry, rank, sparsity, condition-related behaviour, and differences from another finite-element implementation. For nonlinear static analysis, the relevant stiffness is the tangent operator associated with the state at which the output is produced. In buckling analysis, the elastic stiffness can be exported separately from the geometric stiffness.

Matrix output is a diagnostic and verification feature distinct from the ordinary result Field output used for engineering post-processing.

\begin{keywordparameters}
\key{FILE} & optional & Stiffness-matrix output filename. \\
\end{keywordparameters}

\exampleheading
\begin{dialectexamples}
\begin{femasterdeck}
*LOADCASE, TYPE=LINEARSTATIC
*REQUESTSTIFFNESS, FILE=K_service.txt
...
*END
\end{femasterdeck}
\begin{noabaqus}
Abaqus matrix-output controls are currently not translated to FEMaster's native matrix diagnostics.
\end{noabaqus}
\end{dialectexamples}

\keywordpage{REQUESTSTGEOM}

\cmd{REQUESTSTGEOM} requests geometric-stiffness output for \val{LINEARBUCKLING}. \key{FILE} optionally specifies the output filename or base name.

Linear buckling uses a stress-dependent geometric stiffness \(K_G\) generated by the reference preload state. Exporting \(K_G\) exposes the destabilizing contribution of the applied preload and, together with \cmd{REQUESTSTIFFNESS}, permits independent inspection of the generalized buckling problem.

The command changes output only; it does not modify the preload or buckling eigenproblem.

\begin{keywordparameters}
\key{FILE} & optional & Geometric-stiffness output filename/base. \\
\end{keywordparameters}

\exampleheading
\begin{dialectexamples}
\begin{femasterdeck}
*LOADCASE, TYPE=LINEARBUCKLING
*REQUESTSTIFFNESS, FILE=K.txt
*REQUESTSTGEOM, FILE=KG.txt
...
*END
\end{femasterdeck}
\begin{noabaqus}
No corresponding geometric-stiffness export card is currently translated from Abaqus input.
\end{noabaqus}
\end{dialectexamples}

\keywordpage{CONSTRAINTSUMMARY}

\cmd{CONSTRAINTSUMMARY} enables diagnostic reporting of the constraint equations used by the active native loadcase.

When several constraint sources act simultaneously, support collectors, RBM suppression, couplings, ties, connectors, and other kinematic relations can all contribute equations. The summary reports their number and grouping.

In an unexpectedly singular or over-constrained model, the constraint groups can expose a duplicate coupling or a support that constrains a reference node twice more directly than the final sparse matrix.

\cmd{CONSTRAINTSUMMARY} does not choose the enforcement method. \cmd{CONSTRAINTMETHOD} separately determines whether the equations are handled by nullspace projection, elimination, or Lagrange multipliers where supported.

\exampleheading
\begin{dialectexamples}
\begin{femasterdeck}
*LOADCASE, TYPE=LINEARSTATIC
*CONSTRAINTSUMMARY
*CONSTRAINTMETHOD, TYPE=NULLSPACE
...
*END
\end{femasterdeck}
\begin{noabaqus}
No Abaqus output-request card maps to this FEMaster constraint diagnostic.
\end{noabaqus}
\end{dialectexamples}

\section{Abaqus output-request keywords}

The supported Abaqus reader accepts four common file/print request keywords: \cmd{NODE FILE}, \cmd{EL FILE}, \cmd{NODE PRINT}, and \cmd{EL PRINT}. Their requested variable names are read so that the deck remains syntactically acceptable, but the cards do not filter or configure FEMaster result output.

Acceptance of an output card does not imply that FEMaster reproduces the corresponding Abaqus output-file semantics. Post-processing is based on FEMaster's documented result files and Fields.

\keywordpage{NODE FILE}

\cmd{NODE FILE} is accepted in supported Abaqus input but does not create an Abaqus-style node file. Requested variables such as \val{U} or \val{RF} do not change which FEMaster result Fields are produced.

\exampleheading
\begin{dialectexamples}
\begin{nofemaster}
Native FEMaster has no \cmd{NODE FILE} input keyword. Nodal results are written by the enabled FEMaster result writers.
\end{nofemaster}
\begin{abaqusdeck}
*NODE FILE
U, RF
\end{abaqusdeck}
\end{dialectexamples}

\keywordpage{EL FILE}

\cmd{EL FILE} is accepted and ignored in the same sense as \cmd{NODE FILE}. It does not create a separate Abaqus-style element file and does not select a subset of FEMaster element Fields.

\exampleheading
\begin{dialectexamples}
\begin{nofemaster}
Native FEMaster has no \cmd{EL FILE} keyword.
\end{nofemaster}
\begin{abaqusdeck}
*EL FILE
S, E
\end{abaqusdeck}
\end{dialectexamples}

\keywordpage{NODE PRINT}

\cmd{NODE PRINT} is accepted so that imported Abaqus decks may retain common print requests, but it does not generate an Abaqus-formatted nodal print table.

\exampleheading
\begin{dialectexamples}
\begin{nofemaster}
No native \cmd{NODE PRINT} keyword is used for FEMaster result selection.
\end{nofemaster}
\begin{abaqusdeck}
*NODE PRINT
U, RF
\end{abaqusdeck}
\end{dialectexamples}

\keywordpage{EL PRINT}

\cmd{EL PRINT} is the element counterpart of \cmd{NODE PRINT}. The keyword is accepted but does not configure FEMaster result output.

\exampleheading
\begin{dialectexamples}
\begin{nofemaster}
No native \cmd{EL PRINT} keyword is used for FEMaster result selection.
\end{nofemaster}
\begin{abaqusdeck}
*EL PRINT
S, E
\end{abaqusdeck}
\end{dialectexamples}

\section{Native \file{.res} files}

The native \file{.res} file is organized into loadcase and Field blocks. A loadcase header has the form
\begin{syntax}
LC <id>
\end{syntax}
and is followed by the Fields produced by that analysis.

\subsection*{Indexed Field storage}

All standard result Fields associated with model entities are written in indexed form. Every row therefore contains one or more explicit index columns followed by the Field values. The header states both counts:

\begin{syntax}
FIELD, NAME=<name>, TYPE=<type>,
INDEX_COLS=<nindex>, VALUE_COLS=<nvalue>, ROWS=<rows>
<index-1> ... <index-n> <value-1> ... <value-m>
...
END FIELD
\end{syntax}

The indices identify the corresponding model entity rather than an implicit row number. Results therefore remain traceable to the user-visible node and element labels even when the model contains Parts and multiple Instances.

Entities belonging to the implicit default Instance retain their unqualified local identifier, for example \val{8}. Entities belonging to an explicitly named Instance use a qualified identifier of the form

\begin{syntax}
<instance>.<local-id>
\end{syntax}

such as \val{TRUSS_A.8} or \val{TRUSS_B.10}. The same convention is used for node and element identifiers.

The number and meaning of the index columns depend on the Field domain:

\begin{center}
\begin{tabular}{@{}p{0.27\linewidth}p{0.18\linewidth}p{0.45\linewidth}@{}}
\toprule
\textbf{Field type} & \textbf{Index columns} & \textbf{Meaning} \\
\midrule
\val{NODE}          & 1 & node identifier \\
\val{ELEMENT}       & 1 & element identifier \\
\val{ELEMENT_NODAL} & 2 & element identifier, local node index \\
\val{ELEMENT_IP}    & 2 & element identifier, local integration-point index \\
\val{ELEMENT_MP}    & 2 & element identifier, local material-point index \\
\bottomrule
\end{tabular}
\end{center}

Local node, integration-point, and material-point indices are zero based. User node and element identifiers are non-negative labels and may themselves also start at \val{0}; these are separate concepts.

A typical nodal Field therefore looks like:

\begin{syntax}
FIELD, NAME=DISPLACEMENT, TYPE=NODE,
INDEX_COLS=1, VALUE_COLS=6, ROWS=4
               1    0.000000e+00    0.000000e+00 ...
      TRUSS_A.8    4.761905e-02   -4.761905e-02 ...
       TRUSS_B.1    0.000000e+00    0.000000e+00 ...
      TRUSS_B.10    4.761905e-02    0.000000e+00 ...
END FIELD
\end{syntax}

For \val{ELEMENT_NODAL}, \val{ELEMENT_IP}, and \val{ELEMENT_MP}, the first index is the semantic element identifier and the second is the zero-based local position within that element. This keeps the external result format independent of any internal element ordering.

\subsection*{Common result quantities}

Typical Fields depend on analysis type:
\begin{longtable}{@{}p{0.27\linewidth}p{0.63\linewidth}@{}}
\toprule
\textbf{Analysis} & \textbf{Typical results} \\
\midrule
\endhead
Linear static & Displacement, strain, stress, shell top/bottom stresses and resultants where applicable, external forces, reactions, section forces, and related recovered quantities. \\
Nonlinear static & Accepted-increment displacement, stress, reaction-force, and load-factor states plus final displacement, strain, stress, internal/external forces, and reactions. \\
Eigenfrequency & Mode shapes, eigenvalues, eigenfrequencies/frequencies, and participation-related quantities where available. \\
Linear buckling & Buckling mode shapes and buckling factors. \\
Linear transient & Displacement, velocity, and acceleration at written time frames. \\
Linear harmonic & Real and imaginary displacement, stress, and strain amplitudes for the requested excitation frequencies. \\
\bottomrule
\end{longtable}

Generalized nodal motion and force vectors use the six-component nodal order \dofvec. Stress and strain tensors instead use the engineering order
\[
(x,y,z,yz,zx,xy).
\]
The component meaning of raw result data therefore depends on whether the Field contains generalized nodal quantities or stress/strain tensors.

\subsection*{Structural stress and strain recovery}

For solid elements, FEMaster evaluates constitutive stress and strain at the element's actual material integration points. Result recovery extrapolates those values to the element's local nodes with a topology-specific interpolation/extrapolation matrix; the global nodal result is then assembled from the contributions of incident elements. A reduced-integration element can only recover the spatial variation represented by its integration-point data (for example, a single-point \val{C3D8R} supplies a constant recovered value per element). Extrapolated element-nodal values and averaged nodal values therefore need not equal the original integration-point values.

The distinction is particularly important at a material interface, a stress discontinuity, a sharp corner, or a strongly distorted mesh: nodal averaging can smooth physically different element-side values, and extrapolation may produce extrema outside the integration-point range. When comparing results with another solver or experimental stress measures, compare the same stress measure and recovery location, not merely a field with the same name. Local \val{ELEMENT\_IP}, \val{ELEMENT\_NODAL}, and global \val{NODE} domains retain different indexing and recovery semantics in the native \file{.res} format.

FEMR version 1 stores \val{NODE} and \val{ELEMENT} result fields but does not serialize \val{ELEMENT\_NODAL}, \val{ELEMENT\_IP}, or \val{ELEMENT\_MP} fields. The writer skips unsupported domains with a warning; use \file{.res} when the local-field domain is needed. The absence of a local field in FEMR does not imply that FEMaster did not evaluate that quantity internally.

\subsection*{Harmonic results}

A complex harmonic quantity is written as separate real and imaginary components:
\[
\hat q=q_\mathrm{real}+i q_\mathrm{imag}.
\]
Amplitude and phase can be reconstructed from
\[
|\hat q|=\sqrt{q_\mathrm{real}^2+q_\mathrm{imag}^2},
\qquad
\varphi=\operatorname{atan2}(q_\mathrm{imag},q_\mathrm{real}).
\]
Harmonic results distinguish the real part at a chosen phase, the complex amplitude magnitude, and signed components. These quantities are not interchangeable.
Each excitation-frequency sample retains its own six real/imaginary displacement, stress, and strain Fields. An ascending, one-based sweep index is appended to the native Field names (for example, \val{DISPLACEMENT\_REAL\_1} and \val{DISPLACEMENT\_IMAG\_1}); the physical excitation frequency remains the result-frame value. This also keeps samples with repeated frequency values distinct in \file{.res} output.
The node-wise magnitude of the complex displacement vector can be reconstructed from the first three translational components of the paired real and imaginary Fields. The console's response norm instead combines all active degrees of freedom of the full displacement vectors, including rotational degrees of freedom where present, and is not the maximum node-wise displacement amplitude.

The FRD writer preserves this separation rather than discarding the imaginary component. Harmonic displacement Fields are written as \val{DISPREAL} and \val{DISPIMAG}; harmonic stress as \val{STRREAL} and \val{STRIMAG}; and harmonic strain as \val{STRNREAL} and \val{STRNIMAG}. This permits direct visualization or reconstruction of the complex response from the paired FRD blocks.

\section{FEMR binary result format}
\label{sec:femr-format}

FEMR version 1 is FEMaster's seekable binary result container. Multi-byte integers and floating-point values are little-endian. Unknown chunks are skipped using their stored size, allowing compatible format extensions.

Binary output is selected explicitly, for example:
\begin{verbatim}
FEMaster model.inp --output-format femr
FEMaster model.inp --output-format res frd femr
\end{verbatim}
FEMR version 1 always compresses each independently loadable field-data chunk with LZ4. There is no compression setting in the command-line interface.

\subsection{Chunk envelope}

Every chunk begins with a 32-byte envelope containing, in order:
\begin{itemize}
  \item a four-byte chunk code;
  \item a \texttt{uint32} compression code (0 for none, 1 for an LZ4 block);
  \item \texttt{uint64} stored and uncompressed payload sizes;
  \item a \texttt{uint32} CRC32 of the uncompressed payload; and
  \item a reserved \texttt{uint32}, which is zero in version 1.
\end{itemize}
Only \texttt{FDAT} is compressed in version 1. Each field is a separate LZ4 block, so loading one field never decompresses another field.

\subsection{Chunks}

\begin{description}
  \item[HEAD] Must be first. It stores the \texttt{FEMR} magic, major and minor version, byte order, scalar width, default compression and reserved data.
  \item[INST] Stores each instance identifier once together with its length-prefixed UTF-8 name.
  \item[MESH] Starts with \texttt{uint64} node and element counts. Each node stores its dense/global, instance and part-local identifiers as three \texttt{int32} values, followed by three \texttt{float64} coordinates. Each element likewise stores its dense/global, instance and part-local identifiers, followed by a length-prefixed UTF-8 type name, node count and dense/global node identifiers.
  \item[LCAS] Stores the load-case identifier and step type.
  \item[FRAM] Stores load-case identifier, dense frame identifier and value. Eigenfrequency and buckling frame identifiers are derived from the numeric field-name suffix, so equal spectral values remain distinct frames.
  \item[FMET] Stores field identifier, load-case and frame identifiers, name, domain, scalar width, row count and component count.
  \item[FDAT] Stores the field identifier followed by its dense row-major scalar array.
  \item[CSUM] Must be last. Its payload is the CRC32 of every preceding chunk envelope and stored payload byte.
\end{description}

Strings use a \texttt{uint16} byte length followed by UTF-8 bytes. Version 1 supports only the field domains \texttt{UNKNOWN} (0), \texttt{NODE} (1), and \texttt{ELEMENT} (2). Writers skip fields from other domains with a warning; readers reject unsupported domain codes found in a version 1 file.

\subsection{Stable mesh and field indexing}

MESH ordering defines the stable dense result indexing. Node entry $i$ maps to row $i$ of every NODE field, and element entry $i$ maps to row $i$ of every ELEMENT field. Node and element identifiers are nevertheless stored explicitly; they must not be reconstructed from element connectivity. The accompanying instance and part-local identifiers preserve semantic labels such as \texttt{17} and \texttt{bolt.17}.

\subsection{Lazy access}

A reader scans only chunk envelopes and the small metadata chunks. It seeks over MESH and FDAT payloads and retains their offsets. Mesh or field payloads are read, decompressed and checked only when requested.

\begin{verbatim}
import femaster_api as femaster

with femaster.open_results("model.femr") as results:
    stress = results.loadcase(1).frame(20).field("STRESS")
    mesh = results.mesh
\end{verbatim}

\section{FRD visualization output}

FRD provides a CalculiX/CGX-compatible representation of the mesh and supported nodal results. It is intended primarily for visualization. Data that do not naturally map to FRD nodal blocks---for example some element-local, integration-point, material-point, scalar-history, or participation quantities---remain available in the native result representation.

Typical mappings include:
\begin{longtable}{@{}p{0.30\linewidth}p{0.22\linewidth}p{0.38\linewidth}@{}}
\toprule
\textbf{FEMaster result} & \textbf{FRD block} & \textbf{Content} \\
\midrule
\endhead
Displacement / mode / buckling vector & \val{DISP} & Translational and supported generalized displacement components. \\
Harmonic displacement, real / imaginary & \val{DISPREAL} / \val{DISPIMAG} & Separate components of the complex harmonic displacement. \\
Velocity & \val{VELO} & Nodal velocity components. \\
Acceleration & \val{ACCE} & Nodal acceleration components. \\
Reaction force & \val{FORC} & Generalized reaction-force components; nonlinear static writes them for every accepted increment and for the final result frame. \\
Stress & \val{STRESS} & Nodal recovered stress components. \\
Harmonic stress, real / imaginary & \val{STRREAL} / \val{STRIMAG} & Separate components of the complex harmonic stress. \\
Strain & \val{TOSTRAIN} & Nodal recovered strain components. \\
Harmonic strain, real / imaginary & \val{STRNREAL} / \val{STRNIMAG} & Separate components of the complex harmonic strain. \\
Shell resultants & \val{SHR} & Membrane, bending, and transverse-shear resultants. \\
\bottomrule
\end{longtable}

FRD requires numeric node identifiers and therefore cannot use the qualified string identifiers employed by \file{.res}. FEMaster assigns every written node the deterministic FRD identifier
\[
n_{\mathrm{FRD}} = 10^8\,i_{\mathrm{instance}} + n_{\mathrm{local}},
\]
where \(i_{\mathrm{instance}}\) is the numeric Instance identifier and \(n_{\mathrm{local}}\) is the original Part-local node ID. The local node ID must be smaller than \(10^8\). The same FRD node numbering is used consistently for coordinates, element connectivity, and nodal result blocks. Thus identifiers such as \val{100000001} and \val{200000001} represent the same local node label in different Instances; they are output identifiers and do not replace the semantic \val{INSTANCE.local-id} notation used in \file{.res}.

User node and element IDs in FEMaster may include \val{0}. FRD numbering is an output-format concern and does not impose a positive-ID restriction on the FEMaster input model.

Higher-order element connectivity is reordered where necessary to satisfy FRD conventions. For shell resultants, an explicitly assigned Section orientation defines the local output basis. Without an explicit orientation, FEMaster uses a deterministic tangent basis derived from the shell geometry.

Models containing nodes and nodal features but no supported structural elements remain valid FRD models. In this case FEMaster writes the nodal geometry and applicable nodal result fields and simply omits the element-connectivity block.

Transient FRD frames use physical time. Eigenfrequency and buckling frames use the corresponding spectral value where supported by the writer.
